\chapter{Balloon Enteroscopy}

\section*{Overview}
Balloon enteroscopy is an endoscopic examination of the small intestine using an overtube and one or two balloons to advance the scope beyond the reach of standard endoscopy. The exact approach, setting, and preparation depend on the indication, anatomy, and the patient's health.

\section*{Why It Is Done}
It evaluates obscure gastrointestinal bleeding, suspected small-bowel tumors, Crohn disease, abnormal imaging, and lesions requiring biopsy or treatment.

\section*{How It Is Performed}
Under sedation or anesthesia, the scope is introduced through the mouth or anus. Balloon inflation anchors the bowel while the scope and overtube are advanced. Biopsy, polyp removal, dilation, or bleeding treatment can be performed.

\section*{Risks and Recovery}
Pancreatitis, bleeding, perforation, aspiration, infection, and sedation complications are possible. Patients are observed until alert and should follow instructions about eating, driving, and warning symptoms.

\section*{Outcome and Follow-up}
The expected outcome depends on the reason for the procedure, the person's health, and whether additional treatment is needed. The clinician reviews the result with the patient, explains any next steps, and sets follow-up according to the findings and recovery.

\section*{Contraindications and Precautions}
The endoscopy team reviews bowel preparation, sedation risk, cardiopulmonary stability, bleeding medicines, and the ability to complete the examination safely. Suspected perforation or severe instability may require an alternative or urgent specialist management. Anticoagulant and antiplatelet changes require a clinician-directed plan.

\section*{When to Seek Medical Care}
Follow the procedure team's specific instructions. Seek urgent medical assessment for severe or worsening pain, uncontrolled bleeding, fever or signs of infection, trouble breathing, chest pain, fainting, new weakness or confusion, inability to urinate, or any rapidly worsening symptom. The appropriate warning signs depend on the procedure and the patient's medical condition.

\section*{References}
\begin{enumerate}
  \item MedlinePlus. Enteroscopy. U.S. National Library of Medicine; 2025.
  \item Current Medical Diagnosis and Treatment. McGraw-Hill; 2025.
\end{enumerate}
\clearpage
